\section{Fundamentos técnicos del marco RECD}

\subsection{El Reloj Extramental Discreto (RECD)}

El RECD postula que el tiempo extramental no es un continuo newtoniano ni una sucesión de instantes puntuales, sino una secuencia discreta de \textit{articulaciones ordinales} emergentes de la dinámica de persistencia de los sistemas \citep{padilla_recd_preprint_2025,padilla_tau_sistemic_2025}. Estas articulaciones no se definen por el contenido semántico de los eventos, sino por la estructura ordinal de sus trayectorias en el espacio de fases.

La métrica central es el \textbf{Tau Sistémico} $\tau_s$, que cuantifica la persistencia local de una serie o componente en una ventana deslizante. Su valor se deriva de la entropía de permutación (Bandt-Pompe) invertida y escalada de modo que valores altos de $\tau_s$ indican mayor orden y persistencia en la estructura ordinal de la serie.

\subsection{Constructos centrales}

\textbf{Hiper-persistencia.} Un módulo (componente o región) se encuentra en hiper-persistencia cuando su valor actual de $\tau_s$ excede significativamente el promedio de una ventana reciente (típicamente 20 pasos), normalizado por la desviación estándar de dicha ventana. Formalmente:
\[
\text{hyper}_i(k) = \frac{\tau_s^i(k) - \mu_{\text{reciente}}}{\sigma_{\text{reciente}} + \epsilon} > \theta,
\]
donde $\theta = 1.0$ en la configuración de equilibrio utilizada.

\textbf{Trapping estructural (RQA Config B).} Se evalúa mediante análisis de recurrencia cuantitativa aplicado a la serie de $\tau_s$. Dos métricas principales se emplean: Laminaridad (LAM) y Trapping Time (TT). La configuración B (equilibrio) utiliza los umbrales LAM $\geq 0.80$ y TT $\geq 3.4$, junto con aceleraciones de ambas métricas. Un pico se considera \textit{hiperestructural} cuando combina hiper-persistencia con al menos uno de los criterios de trapping estructural.

\textbf{Intensidad compuesta $I$.} Para cada módulo $i$ en el instante $k$ se define
\[
I_i(k) = \text{Peso}_i(k) \times \rho_i(k) \times \Gamma_i(k),
\]
donde Peso corresponde a $\tau_s$, $\rho$ mide la conectividad espacial media con otros módulos y $\Gamma$ captura el exceso de persistencia respecto a la mediana reciente. Los picos estructurales se detectan cuando $I_i$ supera simultáneamente umbrales percentilares en sus tres componentes y se satisface el criterio de hiper-persistencia más trapping.

\textbf{Antisincronización.} Se define como la desviación estándar de los valores de $\tau_s$ entre todos los módulos en un instante dado:
\[
A(k) = \sqrt{\frac{1}{n}\sum_{i=1}^n \bigl(\tau_s^i(k) - \bar{\tau}_s(k)\bigr)^2}.
\]
Valores altos de $A(k)$ indican que los módulos sostienen escalas de persistencia notablemente diferentes (mayor autonomía, mayor fragmentación). Valores bajos indican alineación de escalas (menor fragmentación, mayor permeabilidad entre presentes locales).

\subsection{Parámetros consistentes}

Todos los análisis empíricos que sustentan esta ontología se realizaron con la siguiente configuración invariante:
\begin{itemize}
    \item Ventana para $\tau_s$: 85 pasos.
    \item Ventana de coherencia espacial: 90 pasos.
    \item Umbral de hiper-persistencia: 1.0.
    \item RQA Config B (LAM $\geq 0.80$, TT $\geq 3.4$, aceleraciones correspondientes).
    \item Umbral de percentil para picos estructurales: 94.° percentil con distancia mínima de 35 pasos.
\end{itemize}

Esta constancia paramétrica garantiza que las diferencias observadas entre regímenes y niveles de ruido sean atribuibles a la dinámica del sistema y no a variaciones metodológicas arbitrarias \citep{padilla_rqa_integration_2026,padilla_feigenbaum_theorem_2026}.